\documentclass[12pt,letterpaper]{article}
\usepackage{amsmath}
\usepackage{amsthm}
\usepackage{amsfonts}
\usepackage{amssymb}
\usepackage{amscd}
\usepackage{enumerate}
\usepackage{fancyhdr}
\usepackage{mathrsfs}
\usepackage{bbm}
\usepackage{framed}
\usepackage{mdframed}
\usepackage{cancel}
\usepackage{mathtools}
\usepackage{verbatim}
\usepackage{hyperref}
\usepackage[letterpaper,voffset=-.5in,bmargin=3cm,footskip=1cm]{geometry}
\setlength{\parindent}{0.0in}
\setlength{\parskip}{0.1in}
\allowdisplaybreaks
\headheight 15pt
\headsep 10pt
\newcommand\N{\mathbb N}
\newcommand\Z{\mathbb Z}
\newcommand\R{\mathbb R}
\newcommand\Q{\mathbb Q}
\newcommand\lcm{\operatorname{lcm}}
\newcommand\setbuilder[2]{\ensuremath{\left\{#1\;\middle|\;#2\right\}}}
\newcommand\E{\operatorname{E}}
\newcommand\V{\operatorname{V}}
\newcommand\Pow{\ensuremath{\operatorname{\mathcal{P}}}}

\DeclarePairedDelimiter\ceil{\lceil}{\rceil}
\DeclarePairedDelimiter\floor{\lfloor}{\rfloor}
\newcommand\hint[1]{\textbf{Hint}: #1}
\newcommand\note[1]{\textbf{Note}: #1}
\newcounter{enum22i}
\newenvironment{22enumerate}{\begin{enumerate}[a.]\itemsep0em\setcounter{enumi}{\value{enum22i}}}{\setcounter{enum22i}{\value{enumi}}\end{enumerate}}
\newenvironment{22itemize}{\begin{itemize}\itemsep0em}{\end{itemize}}
\fancypagestyle{firstpagestyle} {
  \renewcommand{\headrulewidth}{0pt}
  \lhead{\textbf{CSCI 0220}}
  \chead{\textbf{Discrete Structures and Probability}}
  \rhead{M. Littman}
}

\fancypagestyle{fancyplain} {
  \renewcommand{\headrulewidth}{0pt}
  \lhead{\textbf{CSCI 0220}}
  \chead{Homework 8}
  \rhead{\textit{ }}
}
\pagestyle{fancyplain}


\begin{document}
  \thispagestyle{firstpagestyle}
  \begin{center}
    {\huge \textbf{Homework 8}}

    {\large Due: \textit{Monday, July 19th at 11:59pm EDT}}
  \end{center}


    All homeworks are due the Monday after their release at 11:59pm EDT on Gradescope.

    Please do not include any identifying information about yourself in the handin, including your Banner ID.

    Be sure to fully explain your reasoning and show all work for full credit. We recommend checking out the sample proofs on the course website to see the level of rigor for which we're looking.
   

\subsection*{Problem 1}
{
\nopagebreak 

{\setcounter{enum22i}{0}

      In this problem, you will prove the main steps of an alternate proof of Fermat's Little Theorem. We will work with an equivalent version of Fermat's Little Theorem, which states that for all integers $a$ and prime numbers $p$, $a^{p} \equiv a \pmod{p}$.
      
      %Note that the version of Fermat's Little Theorem given in class is equivalent, as shown in the steps below.
      %\[a^{p-1} \equiv 1 \pmod{p}\;\;\;\Rightarrow\;\;\;a^{p-1} \cdot a \equiv 1 \cdot a \pmod{p}\;\;\;\Rightarrow\;\;\;a^{p} \equiv a \pmod{p}\]
      
      \begin{22enumerate}
        \item 
          Prove that, for all prime numbers $p$ and all integers 
          $0 < i < p$, $p$ divides ${p \choose i}$.
        \item 
          Prove that, for integers $a$ and $b$ and a prime number $p$, 
          $\displaystyle{(a + b)^{p} \equiv a^{p} + b^{p} \pmod{p}}$.
        \item 
          Complete the proof of Fermat's Little Theorem for positive integers using induction on $a$.
      \end{22enumerate}
    
}
}
\subsection*{Problem 2}
{
{\setcounter{enum22i}{0}
	
Let $p_1$, $p_2$, and $p_3$ be distinct primes and let $r_1$, $r_2$, and $r_3$ be positive integers.
       \begin{22enumerate}   
	    \item Let $n=p_1^{r_1}p_2^{r_2}$.
	    \begin{enumerate}[i.]
	    	\item How many integers between 1 and $n$ (inclusive) are divisible by $p_1$?
	    	\item How many integers between 1 and $n$ (inclusive) are divisible by $p_2$?
	    	\item How many integers between 1 and $n$ (inclusive) are divisible by both $p_2$ and $p_1$?
	    	\item Hence, how many integers between 1 and $n$ (inclusive) are divisible by either $p_1$ or $p_2$ (or both)? Explain the reasoning behind your answer.
	    \end{enumerate}
	     \newpage
      \item Consider a positive integer $n$ whose prime factorization features three distinct primes such that $n=p_1^{r_1}p_2^{r_2}p_3^{r_3}$. How many integers between 1 and $n$ (inclusive) are relatively prime to $n$? Explain the reasoning behind your answer.
      
      \item Hence, what is the value of $\varphi(450)$? Explain the reasoning behind your answer.\\
      \note{Recall that $\varphi$ is the symbol for Euler's totient function.}
      \end{22enumerate} 
	

}

}

\subsection*{Problem 3}
{

{\setcounter{enum22i}{0}

Consider again Saturn's 82 moons. Each of them has some non-negative, integer number of craters. Prove that it is possible to choose 9 of them whose total number of spots is divisible
by 9.\\
\hint{Assign each moon to its number of craters' congruence class and use the Pigeonhole Principle.}
    
}
}

\subsection*{Mind Bender (Extra Credit)}
{
{\setcounter{enum22i}{0}

  Define a set of positive integers to be \textit{yummy} if it
  contains its own cardinality. A set is \textit{minimally yummy} if it is yummy and it has no proper subset which is also yummy.
  Count the number of minimally yummy subsets of $\{1,2, \dots, n\}$.

  \note{A proper subset of a set $S$ is a subset of $S$ that is missing at least one member of $S$.}
    
}
}
  
\end{document}
